%%================================================
%% Abstract - Thai (ภาษาวิชาการ)
%%================================================

โครงงานนี้เสนอแพลตฟอร์ม WheelSense ระบบต้นแบบสภาพแวดล้อมอัจฉริยะสำหรับผู้ใช้เก้าอี้รถเข็นในสถานดูแลผู้สูงอายุ
โดยบูรณาการการรับรู้การเคลื่อนไหว การระบุตำแหน่งภายในอาคาร การบริหารงานเชิงปฏิบัติการ
และการสนับสนุนการตัดสินใจของผู้ดูแล เข้าเป็นสายงานข้อมูลครบวงจร
ภายใต้โครงสร้างพื้นฐานที่ติดตั้งและดูแลรักษาได้จริง
เพื่อแก้ปัญหาความปลอดภัย การสื่อสารข้ามกะ
และการบันทึกข้อมูลแบบกระดาษในบ้านพักคนชราของประเทศไทย

สถาปัตยกรรมของระบบประกอบด้วยสี่ชั้นที่ทำงานประสานกัน
ได้แก่ ชั้นอุปกรณ์ภาคสนามสำหรับตรวจจับการเคลื่อนที่และวัดสัญญาณชีพพื้นฐาน
ชั้นโหนดภายในอาคารสำหรับสแกนสัญญาณวิทยุ (BLE beacon) และเสริมภาพบริบท
ชั้นเซิร์ฟเวอร์ภายในสถานดูแลสำหรับรับเข้า จัดเก็บ
และให้บริการข้อมูลผ่าน API ที่มีการยืนยันตัวตน
และชั้นส่วนติดต่อผู้ใช้ ประกอบด้วยแดชบอร์ดเว็บตามบทบาทและแอปพลิเคชันบนสมาร์ทโฟน
ระบบระบุตำแหน่งใช้แนวทางลายนิ้วมือสัญญาณวิทยุ (RSSI fingerprinting) ร่วมกับอัลกอริทึมเพื่อนบ้านใกล้สุด (K-Nearest Neighbors) ในระดับโซนและห้อง
ส่วนผู้ช่วยปัญญาประดิษฐ์ (AI) แบบประมวลผลภายในเครื่อง (on-device) ถูกบูรณาการผ่าน Model Context Protocol (MCP)
ภายใต้กระบวนการเสนอ--ยืนยัน--ดำเนินการ
เพื่อให้การกระทำที่เปลี่ยนสถานะของระบบต้องได้รับการยืนยันก่อนเสมอ
และเพื่อรักษาความเป็นส่วนตัวของข้อมูลผู้พักอาศัยภายในสถานดูแล

การประเมินระบบต้นแบบดำเนินการที่มหาวิทยาลัยธรรมศาสตร์ครอบคลุมสี่มิติ
ได้แก่ ความทนทานของการจำแนกตำแหน่งจากชุดข้อมูลภาคสนาม;
ความต่อเนื่องและเวลาแฝงของสายงานข้อมูลครบวงจร
ความสอดคล้องในการเลือกเครื่องมือและเวลาแฝงของผู้ช่วย AI จากชุดสถานการณ์ทดสอบ
และประสบการณ์ผู้ใช้ตามบทบาทจากแบบสอบถามออนไลน์
นอกจากนี้ โครงงานได้นำเสนอระบบต้นแบบต่อผู้บริหารและผู้ดูแลของบ้านพักคนชราวาสนะเวศม์
เพื่อเก็บข้อเสนอแนะเชิงการใช้งานจริง โดยไม่มีการติดตั้งระบบ ณ สถานที่ดังกล่าว

ผลการพัฒนาสะท้อนความเป็นไปได้ของการบูรณาการหลายชั้นเทคโนโลยีเข้าเป็นสายงานเดียว
ภายใต้ข้อจำกัดด้านต้นทุน ความซับซ้อนของการติดตั้ง และความเป็นส่วนตัวของข้อมูล
ประเด็นที่ต้องพัฒนาต่อยอด ได้แก่ การสอบเทียบเวลาระหว่างองค์ประกอบย่อย
การขยายชุดกรณีทดสอบให้สะท้อนสภาพงานจริงมากยิ่งขึ้น
และการปรับปรุงประสบการณ์ใช้งานให้เหมาะสมกับผู้ดูแลในบริบทที่หลากหลาย
ซึ่งเป็นแนวทางสำคัญสำหรับการนำระบบไปประยุกต์ใช้ในที่พักอาศัย โรงพยาบาล
และสถานดูแลผู้สูงอายุของประเทศไทยในระยะถัดไป
